\section{Corollaries and design questions}
\label{sec:implications}

The three main results are analytical. The following items are consequences or engineering questions, not recommendations established by the model.

\subsection{Slashable collateral is not necessarily attacker-attributable collateral}

ELVES's rational-adversary argument makes invalid acceptance prohibitively expensive by charging collateral to the attack path \citep{burdges2024}. In a conventional attack, the parties signing or backing an invalid result can be economically aligned with the attacker.

A shared implementation bug changes that mapping. If otherwise-honest guarantors compute the same wrong result and sign it, a later bad verdict can place \emph{their} signatures in the culprit path. The protocol may destroy collateral, yet the actor that selected the crafted input may not bear that loss.

The distinction is:
\[
\boxed{
\text{slashable collateral}
\neq
\text{attacker-attributable collateral}.
}
\]

This does not prove that exploiting the bug is cost-free. It means that the specific expected-profit argument needs an additional mapping from protocol offender to economic beneficiary when honest common-mode faults are admitted.

The wonky region strengthens the issue: the current Gray Paper does not create ordinary culprit/fault offender records for wonky reports. Whether the staking layer imposes a cost on a wonky report is therefore a first-order question for the protocol authors.

\subsection{Guarantor diversity as a design hypothesis}

Every active core has three assigned guarantors and a guarantee carries two or three credentials. If a single fault domain occupies population share \(r\), the probability that at least two of three randomly assigned guarantors belong to that domain is
\[
\Prb(\ge2)=3r^2(1-r)+r^3.
\]
At \(r=0.2\), this is \(0.104\). Across 341 cores, the expectation is approximately 35.5 core assignments per rotation with at least two guarantors from that domain.

This number does \emph{not} establish that an attacker can route a crafted package to one of those cores. It quantifies the frequency of favorable assignments if exploitation opportunities can be targeted or waited for.

One possible mitigation is to require the credential set for a guarantee to span declared implementation/fault-domain classes. This could prevent a single honest client bug from supplying both signatures. But the idea has unresolved issues:
\begin{itemize}
\item implementation identity may be self-declared and therefore gameable;
\item nominally distinct clients can share deeper components;
\item diversity constraints may reduce liveness when one implementation class is unavailable;
\item an adversarial validator can lie about identity, although doing so may force its own signature onto the report.
\end{itemize}
Hence guarantor diversity is a hypothesis to analyse, not a proposed protocol change in this paper.

\subsection{Fault-domain exposure metrics}

The quantity that matters for an unknown future bug cannot be measured directly: by definition we do not know which validators will be wrong on an input not yet discovered.

What can be measured is \emph{exposure} to shared components. Examples include validator share by:
\begin{itemize}
\item client implementation;
\item execution VM;
\item compiler/toolchain;
\item cryptographic or erasure-coding library;
\item major shared dependencies;
\item copied or common code paths.
\end{itemize}

These are upper bounds or proxies for possible fault-domain concentration, not measurements of \(f_x\) for an unknown bug. A useful operational question is therefore whether the largest identifiable shared-component exposure leaves adequate margin below the attribution and escalation boundaries.

\subsection{Announcement shortfall as an early-warning statistic}

JAM's later tranches respond to unmatched announcements. A validator that is offline before announcing is different: it simply reduces the number of initial auditors that appear.

Aggregate announcement counts per report can therefore serve as a statistical early-warning signal for loss of available corrective capacity. The expected count is known from validator count and active-core configuration. A sustained shortfall can indicate that the online population is smaller than the nominal population.

This is not a deterministic fault detector. Audit assignment is random, observations are local before information propagates, and the current formal validator-statistics tuple does not directly expose a finalized audit-count metric. The proposal is simply to monitor the discrepancy between expected and observed audit participation.

\subsection{Rubber-stamping remains an open interface}

The Gray Paper notes that auditing activity is difficult to track directly and delegates validator rewards to a staking subsystem; peer voting on impressions of validator effort is described as the available oracle mechanism \citep{wood2026}.

That leaves a classic verifier's-dilemma-style question \citep{luu2015}: how strongly does the reward signal discriminate actual re-execution from a validator that produces the expected positive judgment without paying the full computation cost?

The present paper does not claim that rubber-stamping is selected in JAM. The reward function and effort-observability mechanism are not sufficiently specified to close that game numerically. Candidate approaches such as verifiable execution artifacts, occasional test reports, or other proof-of-effort mechanisms should therefore be evaluated only after the actual staking design is specified.

\subsection{Load-dependent no-shows are a separate future model}

ELVES treats the no-show/fault share used in its committee-size analysis as exogenous. In a real implementation, no-show probability could rise with audit load. If
\[
u=u(L),
\]
and load \(L\) grows with the final committee size, then
\[
 s_f\approx\frac{s_0}{1-Fu(s_f)}
\]
becomes a feedback equation. This could amplify load under stress, but no claim of a JAM overload transition is made without empirical measurements of \(u(L)\). A testnet stress experiment measuring deadline misses against concurrent audit load would be the appropriate next step.
